A confirmação dos valores na última coluna da Tabela 3.1 é feita aplicando a fórmula de estimativa de rotações fornecida no texto e comparando os resultados com os valores tabelados.

\textbf{Fórmula de número de rotações:}
Para um elemento rodante (por exemplo, uma roda de raio $r$ percorrendo distância $s$), o número de rotações completas é:
\[
N_{\rm rot} = \frac{s}{2\pi r}.
\]
Essa expressão é adimensional quando $s$ e $r$ têm a mesma unidade de comprimento:
\[
\left[\frac{s}{2\pi r}\right] = \frac{L}{1 \cdot L} = 1.\quad\checkmark
\]

\textbf{Procedimento de verificação linha a linha:}
\begin{enumerate}
\item Para cada linha da Tabela 3.1, extrair os dados físicos correspondentes (distância percorrida $s$ e raio $r$ ou parâmetro equivalente).
\item Calcular $N_{\rm rot} = s/(2\pi r)$.
\item Comparar o valor calculado com o valor na última coluna da tabela.
\item Verificar que a diferença se deve apenas a arredondamento.
\end{enumerate}

\textbf{Exemplo numérico (valores ilustrativos):}
Se $s = 628\,\mathrm{m}$ e $r = 1\,\mathrm{m}$:
\[
N_{\rm rot} = \frac{628}{2\pi\cdot 1} \approx \frac{628}{6{,}28} = 100\text{ rota\c{c}\~oes}.
\]
Quando os valores recalculados coincidem com a tabela (dentro de arredondamento), o número estimado de rotações fica confirmado.

